% !TEX root = ../main.tex
% chktex-file 46
\chapter{Zusammenfassung}%
\label{sec:conclusion}

Abschließend wird nun zusammengefasst, inwiefern das in~\ref{sec:intro:goals} beschriebene Ziel einer erweiterbaren, parallelisierbaren Wissensgraph-Konstruktion aus Kommunikationsdaten erreicht wurde.
Hierfür werden in~\ref{sec:conclusion:review} die erreichten Teilziele betrachtet.
Anschließend wird in~\ref{sec:conclusion:todo} ein Ausblick auf Möglichkeiten der Weiterentwicklung des Systems gegeben.

\section{Rückblick}%
\label{sec:conclusion:review}

\begin{figure}[h]
	\centering
	\includegraphics[width=0.65\textwidth]{gfx/conclusion/overview.pdf}
	\caption{Grober Aufbau des Konstruktionsverfahrens}\label{fig:conclusion:overview}
\end{figure}
In Kapitel~\ref{sec:text2kg} wurden Lösungen für die drei wesentlichen Teilprobleme Ontologie, Sprachverarbeitung und Grapherweiterung beschrieben.
Die drei Teillösungen bilden zusammen ein Konstruktionsverfahren, welches dem zu Anfang beschriebenen Aufbau entspricht.

\paragraph{Ontologie}
Die Wissensgraphontologie basiert auf Konzeptgraphen.
Im Gegensatz zu den Fakten-Tripel basierten semantischen Netzen, die lediglich beschreiben, welche Beziehungen zwischen Konzepten existieren, kann mittels Konzeptgraphen auch beschrieben werden, welche Beziehungen nicht, oder nur möglicherweise existieren.
Die konstruierten Wissensgraphen haben also eine deutlich höhere Ausdrucksstärke, als z.~B. NELL oder der Google~Knowledge~Graph.
Diese Ausdrucksstärke ist wichtig, um die Komplexität natürlichsprachlicher Aussagen abbilden zu können.

Damit den konstruierten Konzeptgraphen eine Semantik zugeordnet werden kann, muss die Semantik der verwendeten Relationen beschrieben werden.
Es wurde daher ein möglichst minimales Set von Relationen definiert~\tref{sec:text2kg:ontology:pred}, welches zur Beschreibung einer Vielzahl natürlichsprachlicher Aussagen geeignet ist.

Eine der in~\ref{sec:intro:goals} beschriebenen Anforderungen an das Konstruktionsverfahren ist die Erweiterbarkeit, d.~h.\ die prinzipielle Möglichkeit, nicht natürlichsprachliche Eingaben einzufügen.
Konkret bedeutet dies, dass z.~B. ein Bild in denselben Wissensgraph eingefügt werden kann, in den auch Textnachrichten eingefügt werden.
Da die gewählten Konzeptgraph-Prädikate nicht speziell für natürliche Sprache ausgelegt sind, ist dies prinzipiell möglich;
es gibt keine Prädikate, die bestimmte grammatikalische Konstrukte oder Wortarten abbilden.
Lediglich $label$ und $named$ können nur für textuelle Eingaben verwendet werden, die übrigen Prädikate beschreiben allgemeine Eigenschaften von Konzepten und erfordern daher keine bestimmten Eigabearten.
Wie die Integration anderer Eingabearten im Detail aussehen könnte, war jedoch nicht Teil der Aufgabenstellung und wurde daher nicht näher untersucht.

\paragraph{Sprachverarbeitung}
Für die Transformation einer Nachricht in einen Konzeptgraphen wurde das Stanford~CoreNLP Framework mit einer Menge von Transformationsregeln~\tref{sec:text2kg:nlp:transform} kombiniert, welche die von CoreNLP gefundenen Strukturen auf äquivalente Konzeptgraph-Strukturen abbilden.

\paragraph{Grapherweiterung}
Das Einfügen des aus einer Nachricht extrahierten Konzeptgraphen in einen bestehenden Wissensgraphen erfolgt mittels PSL.\@
Im Rahmen dieser Arbeit wurde PSL erstmals im Kontext einer Konzeptgraph-basierten Wissensgraph-Konstruktion eingesetzt.
Es wurde ein Set von PSL-Regeln~\tref{sec:text2kg:psl:inference} vorgestellt, welches den Einsatz in diesem Kontext exemplarisch aufzeigt.

Neben der Erweiterbarkeit, war auch die Parallelisierbarkeit eines der in~\ref{sec:intro:goals} beschriebenen Ziele.
Durch die Wahl von PSL für die Grapherweiterung wurde diese Anforderung prinzipiell erfüllt.
Das für die Inferenz verwendete ADMM-Verfahren lässt sich gut verteilt implementieren.
In~\cite[Kapitel 10]{Boyd2011} wird beschrieben, wie sich ADMM mit verteilten Programmiermodellen, wie z.~B. MapReduce oder Pregel, implementieren lässt.
Aufbauend darauf, wird in~\cite{Lubell-Doughtie2013} eine Hadoop MapReduce-basierte Implementation vorgestellt.
PSL ist prinzipiell also für den Einsatz in Cluster-Umgebungen geeignet.
Die exemplarische verteilte PSL-Implementation \textit{foxPSL}~\cite{Magliacane2015} demonstriert dies.
In der im Rahmen dieser Arbeit erstellten Implementation wurde die PSL-Referenzimplementation verwendet, da diese am besten dokumentiert ist und die Geschwindigkeit für die verwendeten kleinen Datensets ausreichend ist.

Das Ziel ein erweiterbares, parallelisierbares Wissensgraph-Konstruktionsverfahren für Kommunikationsdaten zu konzipieren wurde somit prinzipiell erreicht.
Nichts desto trotz gibt es noch zahlreiche Möglichkeiten das vorgestellte Verfahren weiterzuentwickeln.

\section{Ausblick}%
\label{sec:conclusion:todo}

Im Folgenden wird ein Überblick über Möglichkeiten der Weiterentwickung und offen gebliebene Fragestellungen gegeben.

\paragraph{Verbessern der Testmethode}
Ein Problem bei der Konzeption des Verfahrens war der Umstand, dass es keine Kommunikationsdaten mit Referenzergebnissen gibt, anhand derer die Auswirkungen von Änderungen am Verfahren hätten bewertet werden können.
Die Qualität vieler Entscheidungen, wie z.~B. der Wahl der Regelgewichte, konnte daher nur auf Basis manuell durchgeführter, stichprobenartiger Tests bewertet werden.

Um dieses Problem zu lösen, könnte ein Kommunikationsdaten-Testset entwickelt werden, mit dem sich die Performance eines Wissensgraph-Konstruktionsverfahrens repräsentativ empirisch messen lässt.
Ein solches Datenset ist wichtig, um ein systematisches Weiterentwickeln des Verfahrens zu ermöglichen.
Das in dieser Arbeit verwendete Enron E-Mail Datenset ist potentiell ein guter Ausgangspunkt hierfür.
Zu klären ist, welche Art von Informationen die Referenzergebnisse enthalten und wie eine umfangreiche Menge solcher Referenzergebnisse konstruiert werden kann.

\paragraph{NLP-Phase}
Im vorgestellten Verfahren wird lediglich ein Teil der von CoreNLP extrahierten Informationen in Konzeptgraph-Strukturen übersetzt.
Durch die Integration bislang ungenutzter Abhängigkeitstypen ließe sich die Qualität der Ergebnisse, wie in~\ref{sec:evaluation:quality:results}~Anfrage~2 gezeigt, deutlich verbessern.

Außerdem sollte die verwendete Heuristik zur Übersetzung von natürlichsprachlicher Negation in negative Kontexte verbessert werden.
In~\ref{sec:evaluation:quality:results}~Anfrage~4 wird dies deutlich.
Ein Ansatz hierfür ist z.~B. das Nutzen des CoreNLP Natural Logic Annotators~\cite{MacCartney2007}, welcher u.~a.\ Quantoren, Negationen und die zugehörigen Scopes, d.~h.\ die modifizierten Token, erkennt.

\paragraph{PSL-Phase}
Die verwendeten domänenspezifischen Regeln und das domänenspezifische Vorwissen sind bislang sehr rudimentär.
In künftigen Arbeiten könnten vorhandene Wissensbasen, wie z.~B. NELL~\cite{Carlson2010} oder YAGO~\cite{YAGO}, in die PSL-Inferenz integriert werden.

Neben der Nutzung domänenspezifischen Wissens, besteht außerdem Verbesserungspotential bei der Nutzung von Kontexten in der Inferenz.
Bislang werden Kontexte nur benutzt, um die Richtung von $inst$-Relationen zu ermitteln.
Das Inferieren des $actual$-Attributs von Möglichkeitskontexten~\tref{sec:text2kg:ontology:modal} zur Bestimmung der Glaubwürdigkeit von Aussagen wäre z.~B. eine sinnvolle Erweiterung des Verfahrens.

\paragraph{Inferenz}
Wie~\treft{sec:evaluation:time} gezeigt hat, besteht noch Verbesserungsbedarf bei der Geschwindigkeit der PSL-Inferenz, bevor das Verfahren in der Praxis einsetzbar ist.
Um dies zu erreichen, ist die Kombination verschiedener Ansätze denkbar:
\begin{enumerate}
	\item Nutzen einer Cluster-fähigen ADMM-Implementation, statt der PSL-Referenz\-implementation.
		Hierfür könnte z.~B. die zuvor erwähnte Hadoop MapReduce Implementation~\cite{Lubell-Doughtie2013} oder foxPSL~\cite{Magliacane2015} verwendet werden.
	\item Weiterentwicklung von BOCI, sodass wachsende Eingabemengen unterstützt werden.
		Bereits im ursprünglichen Paper~\cite[Abschnitt 6]{Pujara2015} wird dies als offene Fragestellung für folgende Arbeiten genannt.
		Mittels des Inferenz-Budgets könnte dann die Inferenzdauer nahezu beliebig verringert werden, sofern dafür eine entsprechende Verschlechterung der Inferenzqualität hingenommen wird.
		In bisherigen Tests konnten Inferenz-Laufzeiten durch BOCI um mehr als 65\% reduziert werden, ohne dabei die Inferenzqualität deutlich zu verschlechtern~\cite[Abschnitt 5]{Pujara2015}.
	\item Effizienterer Umgang mit transitiven Relationen.
		Bei der Inferenz werden $inst$-Atome aufgrund der Transitivität von $inst$ instanziiert.
		Während der Inferenz sind diese Atome wichtig, sie werden momentan allerdings auch Teil des Wissensgraphen.
		Das Entfernen dieser redundanten $inst$-Atome, würde die Kantenanzahl des in~\ref{sec:evaluation:time} konstruierten Graphen um mehr als 75\% verringern.
		Während der Inferenz würden die transitiven $inst$-Relationen dann nur temporär instanziiert.
		Inwiefern durch eine derartige Erweiterung der PSL-Inferenz neben der Graph-Größe auch die Inferenzdauer verringert werden kann, ist zu untersuchen.
\end{enumerate}

Wie soeben gezeigt, gibt es viele Möglichkeiten das vorgestellten Wissensgraph-Konstruktionsverfahren weiterzuentwickeln.
Diese Arbeit ist also primär als ein Ausgangspunkt zu verstehen, der zeigt, dass die Kombination von Konzeptgraphen, NLP und PSL prinzipiell für die automatisierte Konstruktion von Wissensgraphen geeignet ist.
Aufbauend auf dieser Grundidee kann nun in künftigen Arbeiten ein erweitertes, praxistaugliches Verfahren entwickelt werden.
